\section{Related work}
\label{sec:related}
\para{Decision-making Agents in Minecraft.}
Minecraft is an open-ended 3D world with incredibly flexible game mechanics supporting a broad spectrum of activities. Built upon notable Minecraft benchmarks~\cite{fan2022minedojo,guss2019minerl,guss2019minerl2,guss2021minerl,kanervisto2022minerl,johnson16malmo}, Minecraft learning algorithms can be divided into two categories: 
1) Low-level controller: Many prior efforts leverage hierarchical reinforcement learning to learn from human demonstrations~\cite{lin2021juewumc,mao2021seihai,skrynnik2021hierarchical}. 
Kanitscheider et al.~\cite{kanitscheider2021multitask} design a curriculum based on success rates, but its objectives are limited to curated items.
MineDojo~\cite{fan2022minedojo} and VPT~\cite{openai2022vpt} utilize YouTube videos for large-scale pre-training. 
DreamerV3~\cite{hafner2023mastering}, on the other hand, learns a world model to explore the environment and collect diamonds.
2) High-level planner: Volum et al.~\cite{volum2022craft} leverage few-shot prompting with Codex~\cite{openai2021codex} to generate executable policies, but they require additional human interaction. Recent works leverage LLMs as a high-level planner in Minecraft by decomposing a high-level task into several subgoals following Minecraft recipes~\cite{wang2023describe,nottingham2023embodied,yuan2023plan4mc}, thus lacking full exploration flexibility.
Like these latter works, \voyager also uses LLMs as a high-level planner by prompting GPT-4 and utilizes Mineflayer~\cite{mineflayer} as a low-level controller following Volum et al.~\cite{volum2022craft}. Unlike prior works, \voyager employs an automatic curriculum that unfolds in a bottom-up manner, driven by curiosity, and therefore enables open-ended exploration.

\para{Large Language Models for Agent Planning.} 
Inspired by the strong emergent capabilities of LLMs, such as zero-shot prompting and complex reasoning~\cite{bommasani2021opportunities,brown2020gpt3,raffel2020t5,wei2022emergent,chowdhery2022palm,chung2022flan}, embodied agent research~\cite{duan2022survey, batra2020rearrangement, ravichandar2020recent, collins2021review} has witnessed a significant increase in the utilization of LLMs for planning purposes.
Recent efforts can be roughly classified into two groups.
1) Large language models for robot learning:
Many prior works apply LLMs to generate subgoals for robot planning~\cite{huang2022language,huang2022language,ahn2022saycan,min2021film,blukis2021a}. 
Inner Monologue~\citep{huang2022inner} incorporates environment feedback for robot planning with LLMs.
Code as Policies~\cite{liang2022code} and ProgPrompt~\cite{singh2022progprompt} directly leverage LLMs to generate executable robot policies. 
VIMA~\cite{jiang2022vima} and PaLM-E~\cite{driess2023palme} fine-tune pre-trained LLMs to support multimodal prompts.
2) Large language models for text agents:
ReAct~\cite{yao2022react} leverages chain-of-thought prompting~\cite{wei2022chain} and generates both reasoning traces and task-specific actions with LLMs. 
Reflexion~\cite{shinn2023reflexion} is built upon ReAct~\cite{yao2022react} with self-reflection to enhance reasoning. 
AutoGPT~\cite{autogpt} is a popular tool that automates NLP tasks by crafting a curriculum of multiple subgoals for completing a high-level goal while incorporating ReAct~\cite{yao2022react}'s reasoning and acting loops. 
DERA~\cite{nair2023dera} frames a task as a dialogue between two GPT-4~\cite{openai2023gpt4} agents. 
Generative Agents~\cite{park2023generative} leverages ChatGPT~\cite{chatgpt} to simulate human behaviors by storing agents' experiences as memories and retrieving those for planning, but its agent actions are not executable. SPRING~\cite{wu2023spring} is a concurrent work that uses GPT-4 to extract game mechanics from game manuals, based on which it answers questions arranged in a directed acyclic graph and predicts the next action.
All these works lack a skill library for developing more complex behaviors, which are crucial components for the success of \voyager in lifelong learning. 


\para{Code Generation with Execution.}
Code generation has been a longstanding challenge in NLP~\cite{openai2021codex,nijkamp2022conversational,le2022coderl,chowdhery2022palm,brown2020gpt3}, with various works leveraging execution results to improve program synthesis. Execution-guided approaches leverage intermediate execution outcomes to guide program search~\cite{chen2019execution,chen2021latent,ellis2019write}. Another line of research utilizes majority voting to choose candidates based on their execution performance~\cite{li2022competitionlevel,cobbe2021training}. Additionally, LEVER~\cite{ni2023lever} trains a verifier to distinguish and reject incorrect programs based on execution results. CLAIRIFY~\cite{skreta2023errors}, on the other hand, generates code for planning chemistry experiments and makes use of a rule-based verifier to iteratively provide error feedback to LLMs. \voyager distinguishes itself from these works by integrating environment feedback, execution errors, and self-verification (to assess task success) into an iterative prompting mechanism for embodied control.
